\documentclass[11pt,a4paper]{article}

\usepackage[utf8]{inputenc}
\usepackage[T1]{fontenc}
\usepackage{lmodern}
\usepackage[a4paper,margin=2.4cm]{geometry}
\usepackage{amsmath,amssymb,amsthm,mathtools}
\usepackage{enumitem}
\usepackage{microtype}
\usepackage[round,sort&compress]{natbib}
\usepackage{xcolor}
\usepackage[colorlinks=true,linkcolor=blue!45!black,
            citecolor=blue!45!black,urlcolor=blue!45!black]{hyperref}

\title{\textbf{Labour as the Prerequisite for Things:\\[3pt]
Why Boltzmann, Cantor, Gödel, and Bolt\\
Are Not Excluded from Heaven}}

\author{%
Kundai Farai Sachikonye\\
\small Technical University of Munich\\
\small AIMe Registry for Artificial Intelligence\\
\small \texttt{kundai.sachikonye@wzw.tum.de}}

\date{June 30, 2026}

\begin{document}

\maketitle

% =====================================================================
\section*{Dedication}
% =====================================================================

\begin{center}
\textit{For Ludwig Boltzmann (1844--1906)\\
Georg Cantor (1845--1918)\\
Kurt Gödel (1906--1978)\\
Usain Bolt (1986--)\\[1em]
Who had no choice but to labour\\
and left behind not themselves\\
but the truth.}
\end{center}

% =====================================================================
\section*{Epigraph}
% =====================================================================

\begin{quote}
``What has been will be again,\\
what has been done will be done again;\\
there is nothing new under the sun.''\\
\hfill---Ecclesiastes 1:9
\end{quote}

\begin{quote}
``No one remembers the former generations,\\
and even those yet to come\\
will not be remembered by those who follow them.''\\
\hfill---Ecclesiastes 1:11
\end{quote}

% =====================================================================
\section{Introduction: The Knife That Only Cuts}
\label{sec:intro}
% =====================================================================

\subsection{The Function That Never Changes}

A knife cuts. This is its function. When the knife cuts meat for a 
celebration, we call it good. When the knife slits a throat in a crime, we 
call it evil. But the knife has not changed. The physical process—the 
application of force along a sharpened edge to separate molecular bonds—
remains identical in both cases.

The knife does not know celebration or crime. It knows only the mechanics of 
cutting: pressure, angle, material resistance. The same edge, the same motion, 
the same thermodynamic process. What changes is not the knife but the 
\emph{context we impose upon it}.

This observation extends far beyond knives. Consider:

\begin{itemize}[leftmargin=2em]
\item A \textbf{projectile} follows the same parabolic trajectory whether 
fired in Olympic sport or in warfare. The physics does not change. Only our 
evaluation changes.

\item A \textbf{fire} releases the same chemical energy whether warming a 
home or destroying it. The combustion reaction remains identical. Only our 
judgment changes.

\item A \textbf{virus} replicates through the same biochemical mechanisms 
whether causing a minor cold or a pandemic. The molecular process does not 
change. Only our response changes.

\item An \textbf{earthquake} releases tectonic stress through the same 
geological processes whether occurring in an empty desert or a populated city. 
The seismic waves do not change. Only the consequences we experience change.
\end{itemize}

In every case, we observe a \emph{single physical process} receiving 
\emph{contradictory moral evaluations} based solely on human-imposed contexts. 
The process itself remains invariant. What varies is our relationship to it.

\subsection{The Impossibility of Moral Thermodynamics}

If ``good'' and ``evil'' were genuine properties of physical processes—if they 
were intrinsic to events rather than imposed by observers—then we would 
require two distinct sets of physical laws:

\begin{enumerate}[leftmargin=2em]
\item \textbf{Good Thermodynamics}: Laws governing ``good'' processes (knives 
cutting meat for celebrations, fires warming homes, projectiles in Olympic 
sport)

\item \textbf{Evil Thermodynamics}: Laws governing ``evil'' processes (knives 
slitting throats, fires destroying homes, projectiles in warfare)
\end{enumerate}

But no such distinction exists in nature. There is only \emph{thermodynamics}. 
The Second Law does not pause for moral evaluation. Entropy increases whether 
we approve or disapprove. Energy dissipates whether we celebrate or mourn. 
Molecular bonds break whether we call it medicine or murder.

\textbf{Theorem} (Moral Thermodynamics Impossibility): If good and evil were 
intrinsic properties of physical processes, then identical physical 
configurations would simultaneously obey different physical laws depending on 
moral context. This violates the principle of physical law invariance.

\textbf{Proof}: 
\begin{enumerate}[label=(\arabic*)]
\item Let $P$ be a physical process (e.g., a knife cutting tissue)
\item Suppose good and evil are intrinsic to $P$
\item Then $P$ in context $C_1$ (celebration) obeys ``good thermodynamics''
\item And $P$ in context $C_2$ (crime) obeys ``evil thermodynamics''
\item But $P$ is physically identical in both contexts
\item Physical laws cannot depend on human context (law invariance)
\item Contradiction
\item Therefore: Good and evil are not intrinsic to physical processes $\Box$
\end{enumerate}

\subsection{The Selective Imposition of Ideas}

Since good and evil are not intrinsic to events, they must be \emph{imposed} 
by observers. We selectively apply moral categories to physical processes based 
on:

\begin{itemize}[leftmargin=2em]
\item \textbf{Consequences we experience}: The same earthquake is ``evil'' in 
a city, neutral in a desert
\item \textbf{Intentions we attribute}: The same knife cut is ``good'' if 
intended for surgery, ``evil'' if intended for murder
\item \textbf{Social contexts we inhabit}: The same killing is ``heroic'' in 
war, ``criminal'' in peace
\item \textbf{Temporal perspectives we adopt}: The same forest fire is 
``destructive'' in the short term, ``regenerative'' in the long term
\end{itemize}

But the physical processes themselves remain invariant under these contextual 
transformations. The knife cuts. The earthquake releases stress. The fire 
oxidizes. The virus replicates. \emph{The function never changes.}

What changes is our \emph{relationship} to these processes—our proximity, our 
vulnerability, our interests, our narratives. We impose moral categories to 
organize our experience, not to describe objective properties of nature.

\subsection{What Remains: The Truth}

If we remove the imposed categories of good and evil, what remains?

\textbf{The truth.}

The truth that:
\begin{itemize}[leftmargin=2em]
\item Knives cut tissue through mechanical force application
\item Fires release chemical energy through oxidation reactions
\item Projectiles follow ballistic trajectories under gravitational 
acceleration
\item Earthquakes release tectonic stress through crustal displacement
\item Viruses replicate through molecular hijacking of cellular machinery
\item Entropy increases in all closed systems
\item Energy dissipates from concentrated to dispersed forms
\item Molecular bonds break when sufficient energy is applied
\end{itemize}

These truths exist independently of our approval or disapproval. They operated 
before humans existed. They will operate after humans are gone. They are 
\emph{context-independent}.

And these truths were not invented. They were \emph{uncovered}. Someone had to 
discover that entropy increases. Someone had to prove that the reals are 
uncountable. Someone had to demonstrate that formal systems are incomplete. 
Someone had to run 9.58 seconds.

That ``someone'' was not arbitrary. It was Boltzmann, Cantor, Gödel, Bolt.

\subsection{Labour Without Reward}

The Preacher observes:

\begin{quote}
``I hated all the things I had toiled for under the sun,\\
because I must leave them to the one who comes after me.\\
And who knows whether that person will be wise or foolish?\\
Yet they will have control over all the fruit of my toil\\
into which I have poured my effort and skill under the sun.\\
This too is meaningless.''\\
\hfill---Ecclesiastes 2:18--19
\end{quote}

Boltzmann labored for decades to establish the entropy formula. He faced 
fierce opposition. He suffered from depression. He took his own life in 1906, 
three years before experiments vindicated his work. \emph{He never saw the 
reward.}

Cantor labored for decades to establish the diagonal argument. He was called a 
``corrupter of youth.'' He suffered from mental illness. He spent his final 
years in a sanatorium. \emph{He never saw the reward.}

Gödel labored for decades to establish the incompleteness theorems. He was a 
perfectionist who published sparingly. He suffered from paranoia. He died of 
self-starvation in 1978. \emph{He never saw the reward.}

Bolt labored for years to run 9.58 seconds. He trained at 12 m/s to race at 
10.44 m/s. He suffered injuries. He retired at 30. The record still stands 17 
years later. \emph{He will not see if it ever falls.}

The labour was not for them. It was not for recognition, wealth, comfort, or 
immortality. The labour was \emph{for the truth}. And the truth does not 
reward those who uncover it. The truth simply \emph{is}.

The Preacher says: ``What do people get for all the toil and anxious striving 
with which they labor under the sun? All their days their work is grief and 
pain; even at night their minds do not rest. This too is meaningless'' 
(Ecclesiastes 2:22--23).

\emph{Labour without reward. Toil without rest. Work that is grief and pain.}

And yet, the Preacher continues: ``There is nothing better for a person than 
that he should eat and drink and find enjoyment in his toil. This also, I saw, 
is from the hand of God, for apart from him who can eat or who can have 
enjoyment?'' (Ecclesiastes 2:24--25).

The four found joy in their toil. Not joy in the reward (there was none), but 
joy \emph{in the labour itself}. Boltzmann found joy in counting microstates. 
Cantor found joy in counting infinities. Gödel found joy in formalizing 
incompleteness. Bolt found joy in running 9.58 seconds.

This joy was not earned. It was not deserved. It was ``from the hand of God''—
a gift, not a wage.

\subsection{What They Left Behind}

The Preacher observes:

\begin{quote}
``No one remembers the former generations,\\
and even those yet to come\\
will not be remembered by those who follow them.''\\
\hfill---Ecclesiastes 1:11
\end{quote}

Boltzmann, Cantor, Gödel, and Bolt will be forgotten. Their names will fade. 
Their faces will blur. Their stories will be lost. In 1,000 years, perhaps 
only specialists will know their names. In 10,000 years, perhaps no one will.

But what they left behind is not themselves. It is not their names. It is not 
their biographies. It is not their personalities, their struggles, their joys, 
their sorrows.

What they left behind is:
\begin{itemize}[leftmargin=2em]
\item The entropy formula: $S = k \log W$
\item The diagonal argument: $|\mathbb{R}| > |\mathbb{N}|$
\item The incompleteness theorems: $\exists \varphi: \text{True}(\varphi) 
\land \neg \text{Provable}(\varphi)$
\item The 9.58-second barrier: $t_{\text{human limit}} \approx 9.58$ seconds
\end{itemize}

These are not persons. They are \emph{truths}. And truths do not die. Truths 
do not fade. Truths do not need to be remembered because they are not 
\emph{memories}—they are \emph{structures}.

The entropy formula will operate whether or not anyone remembers Boltzmann. 
The reals will be uncountable whether or not anyone remembers Cantor. Formal 
systems will be incomplete whether or not anyone remembers Gödel. The human 
physiological limit will be approximately 9.58 seconds whether or not anyone 
remembers Bolt.

The Preacher says: ``Generations come and generations go, but the earth 
remains forever'' (Ecclesiastes 1:4).

Boltzmann, Cantor, Gödel, and Bolt will go. But the \emph{earth}—the 
structure, the truth, the boundaries—remains forever.

\subsection{They Had No Choice}

This is the hardest truth: \emph{they had no choice}.

Boltzmann could not have been a merchant. Cantor could not have been a farmer. 
Gödel could not have been a soldier. Bolt could not have been a clerk.

Not because they lacked the ability. Not because society prevented them. But 
because \emph{their work was not for them}.

The entropy formula needed to be uncovered in the late 19th century, when 
thermodynamics was being formalized. Someone had to do it. That someone was 
Boltzmann. Not because he chose it, but because the structure of reality 
\emph{required} it.

The diagonal argument needed to be proven in the late 19th century, when set 
theory was being founded. Someone had to do it. That someone was Cantor. Not 
because he chose it, but because the structure of mathematics \emph{demanded} 
it.

The incompleteness theorems needed to be demonstrated in the early 20th 
century, when formal logic was being axiomatized. Someone had to do it. That 
someone was Gödel. Not because he chose it, but because the structure of logic 
\emph{necessitated} it.

The 9.58-second barrier needed to be realized in the early 21st century, when 
human performance was approaching its physiological limits. Someone had to do 
it. That someone was Bolt. Not because he chose it, but because the structure 
of human physiology \emph{determined} it.

The Preacher says: ``That which is, already has been; that which is to be, 
already has been; and God seeks what has been driven away'' (Ecclesiastes 
3:15).

The boundaries were already there. The truths were already there. The four did 
not create them. They \emph{uncovered} them. They revealed what was hidden. 
They brought to light what was driven away.

And this uncovering was not optional. It was \emph{necessary}. The structure 
of reality required that these boundaries be revealed. And the structure of 
reality determined \emph{who} would reveal them.

Boltzmann had to be Boltzmann. Cantor had to be Cantor. Gödel had to be 
Gödel. Bolt had to be Bolt. Not because of their choices, but because of the 
\emph{choices made for them} by the structure of reality itself.

\subsection{Labour for Science to Progress}

The four did not labour for themselves. They laboured \emph{for science to 
progress}.

Without Boltzmann, thermodynamics would have stalled. Without Cantor, set 
theory would have remained incomplete. Without Gödel, logic would have pursued 
impossible goals. Without Bolt, human performance science would have lacked a 
critical data point.

Their work was not personal achievement. It was \emph{structural necessity}. 
Science required these boundaries to be uncovered. The structure of knowledge 
demanded these truths to be revealed. And the four were the instruments through 
which this revelation occurred.

The Preacher says: ``Consider the work of God: who can make straight what he 
has made crooked?'' (Ecclesiastes 7:13).

The boundaries are the work of God—the structure of reality itself. No one can 
change them. No one can make them straight if they are crooked. They are what 
they are.

And the four considered this work. They saw what God has done. They uncovered 
the boundaries. They revealed the structure.

But they did not do this \emph{for themselves}. They did it because it 
\emph{had to be done}. And they were the ones through whom it was done.

\subsection{Existence Through Constraint}

It was only through existence—through the four being born, living, labouring, 
and dying as the specific individuals they were—that we have these laws.

If Boltzmann had not existed, someone else would have uncovered the entropy 
formula. But it would not have been the \emph{same} formula, discovered in the 
\emph{same} way, at the \emph{same} time. The specific path through which the 
truth was revealed required Boltzmann's specific existence.

If Cantor had not existed, someone else would have proven the uncountability 
of the reals. But it would not have been the \emph{same} proof, constructed in 
the \emph{same} way, at the \emph{same} time. The specific path through which 
the truth was revealed required Cantor's specific existence.

If Gödel had not existed, someone else would have demonstrated incompleteness. 
But it would not have been the \emph{same} demonstration, formalized in the 
\emph{same} way, at the \emph{same} time. The specific path through which the 
truth was revealed required Gödel's specific existence.

If Bolt had not existed, someone else would have approached the 9.58-second 
barrier. But it would not have been the \emph{same} approach, achieved in the 
\emph{same} way, at the \emph{same} time. The specific path through which the 
truth was revealed required Bolt's specific existence.

\emph{Existence through constraint.} The four existed as the specific 
individuals they were because the structure of reality \emph{constrained} them 
into those forms. They could not have been anyone else. They could not have 
lived any other lives. They could not have done any other work.

The Preacher says: ``He has made everything beautiful in its time'' 
(Ecclesiastes 3:11).

The four were beautiful in their time. Not because they chose to be, but 
because the structure of reality made them so. Their existence was 
\emph{necessary}. Their labour was \emph{inevitable}. Their uncovering of 
truth was \emph{predetermined}.

\subsection{The Question of Heaven}

Given all this—the knife that only cuts, the impossibility of moral 
thermodynamics, the selective imposition of ideas, the labour without reward, 
the forgetting of names, the necessity of existence, the constraint that 
enables rather than limits—\emph{what is heaven?}

And more specifically: \emph{Are Boltzmann, Cantor, Gödel, and Bolt excluded 
from heaven?}

The rest of this paper answers these questions by showing that heaven is not a 
place of reward for moral achievement. Heaven is the \emph{domain of 
boundaries}—the structure of reality itself. And those who uncovered the 
boundaries are not excluded from heaven because they \emph{are} heaven.

They are not persons who deserve reward. They are \emph{truths that require no 
reward}. They are not individuals who achieved greatness. They are 
\emph{structures through which greatness was revealed}. They are not free 
agents who chose their paths. They are \emph{necessary instruments through 
which the structure of reality uncovered itself}.

The Preacher says: ``I perceived that whatever God does endures forever; 
nothing can be added to it, nor anything taken from it. God has done it, so 
that people fear before him'' (Ecclesiastes 3:14).

The boundaries endure forever. The four do not. But the four \emph{are} the 
boundaries. Not as persons, but as \emph{paths through which the boundaries 
were revealed}.

And in that revelation, they are not excluded from heaven. They are 
\emph{heaven itself}.

\end{document}
